\chapter{Customer Relationship Management}
\label{ch:crm}

\chapterguide
  {You should know the sales-order process and the difference between a potential sale and a confirmed order.}
  {explain the purpose of CRM, distinguish operational CRM activities, connect CRM to ERP data, and evaluate how CRM supports acquisition, retention, sales, marketing, and service.}

\begin{sourcebasis}
This chapter follows Tutorial 4 and the CRM section of Chapter 3 in the
textbook. Tutorial 4 uses Luna Electronics and asks students to justify CRM,
propose a CRM system, explain its features, and support the proposal with a case
study.
\end{sourcebasis}

\section{CRM is a relationship process, not just a contact list}

\term{Customer Relationship Management (CRM)} is a set of processes and systems
for managing interactions with prospects and customers across the whole
relationship. A contact database answers ``Who is this customer?'' A CRM system
must also answer what happens next, and who owns it.

\begin{erpfig}[Contact list versus CRM]{The same customer record, two levels of
ambition. Everything in the right-hand column is a question about action; a
contact list cannot answer any of them.}{fig:contact-vs-crm}
\begin{tikzpicture}[x=1cm,y=1cm]
\draw[draw=UTPMuted!45,fill=UTPPaleGray,rounded corners=2.5mm] (0,-0.55) rectangle (6.4,3.6);
\node[erpcap,color=UTPMuted] at (3.2,3.25) {\bfseries Contact list};
\node[erpquiet,minimum width=5.4cm,minimum height=0.7cm] at (3.2,2.5) {Name and address};
\node[erpquiet,minimum width=5.4cm,minimum height=0.7cm] at (3.2,1.65) {Telephone and email};
\node[erpquiet,minimum width=5.4cm,minimum height=0.7cm] at (3.2,0.8) {Contact person};
\node[erpnote,color=UTPMuted] at (3.2,0.0) {answers \emph{who}};

\draw[draw=UTPBlue!45,fill=UTPPaleBlue!45,rounded corners=2.5mm] (7.6,-0.55) rectangle (14.6,3.6);
\node[erpcap,color=UTPBlue] at (11.1,3.25) {\bfseries CRM};
\foreach \i/\t in {0/{Which opportunities are open, and worth how much?},
                   1/{Who owns the next action, and by when?},
                   2/{What has this customer bought, been promised, complained about?}}{
  \pgfmathsetmacro{\y}{2.5-\i*0.85}
  \node[erpstep,minimum width=6.4cm,text width=6.0cm,minimum height=0.7cm] at (11.1,\y) {\t};}
\node[erpnote,color=UTPBlue] at (11.1,0.0) {answers \emph{what next}};
\draw[erpflow] (6.6,1.5) -- (7.4,1.5);
\end{tikzpicture}
\end{erpfig}

\begin{intuition}
CRM gives the organisation a ``single face to the customer.'' The customer should
not receive three contradictory prices from Sales, Marketing, and an executive
simply because those people do not share interaction history.
\end{intuition}

\section{Core CRM activities}

\begin{erptable}
\begin{tabularx}{\linewidth}{@{}>{\bfseries\leavevmode\color{ERPNavy}}p{0.25\linewidth}X@{}}
\toprule
\rowcolor{ERPPaleBlue!92}
CRM activity & Purpose \\
\midrule
One-to-one marketing & Tailor offers and communication to a customer's profile or history; supports cross-selling and up-selling. \\
Sales force automation & Record contacts, assign leads, track opportunities, manage follow-ups, and forecast sales activity. \\
Campaign management & Organise marketing campaigns and analyse their outcomes. \\
Marketing content / knowledge & Give staff consistent product and promotional information. \\
Service and call-centre support & Provide representatives with customer and product history so issues can be resolved without starting from zero. \\
\bottomrule
\end{tabularx}
\end{erptable}

\section{Lead, opportunity, quotation, and order are different states}

A common beginner mistake is to treat every prospect as a sales order. CRM
exists partly because sales work begins long before a customer commits --- and
because most prospects never do.

\begin{erpfig}[The sales funnel]{Each stage is a smaller, better-qualified set.
The illustrative counts show why a pipeline cannot be read as revenue: only the
last two stages are commitments.}{fig:sales-funnel}
\begin{tikzpicture}[x=1cm,y=1cm]
\foreach \i/\w/\lbl/\n/\c in {%
  0/11.4/{Lead --- an expression of interest}/{200}/UTPMuted,
  1/9.4/{Qualified opportunity --- actively pursued}/{60}/UTPBlue,
  2/7.4/{Quotation --- a commercial offer}/{25}/UTPBlue,
  3/5.4/{Sales order --- customer commitment}/{12}/UTPGreen,
  4/5.4/{Ongoing relationship --- service, repeat, referral}/{12}/UTPGreen}{
  \pgfmathsetmacro{\y}{-\i*1.02}
  \node[draw=\c,fill=\c!12,rounded corners=2mm,minimum width=\w cm,
        minimum height=0.8cm,font=\scriptsize\bfseries\color{UTPNavy}] (f\i) at (6.2,\y) {\lbl};
  \node[erpchip,anchor=west] at (12.8,\y) {\n\ deals};
}
\foreach \i in {0,...,3}{\pgfmathtruncatemacro{\j}{\i+1}\draw[erpflow] (f\i)--(f\j);}
\draw[decorate,decoration={brace,amplitude=4pt,mirror},draw=UTPMuted,semithick]
  (0.0,0.42) -- (0.0,-2.5);
\node[erpcap,anchor=east,text width=1.5cm,align=right] at (-0.25,-1.0) {\bfseries CRM\\ owns these};
\draw[decorate,decoration={brace,amplitude=4pt,mirror},draw=UTPGreen,semithick]
  (0.0,-2.65) -- (0.0,-4.5);
\node[erpcap,anchor=east,text width=1.5cm,align=right,color=UTPGreen] at (-0.25,-3.6) {\bfseries Sales\\ and ERP};
\end{tikzpicture}
\end{erpfig}

A \term{lead} is a potential contact. An \term{opportunity} is a deal being
actively pursued. A \term{quotation} is a commercial offer. A \term{sales order}
is a confirmed commitment. After the sale, service and marketing continue the
relationship.

\section{Why CRM becomes more valuable when integrated with ERP}

CRM holds customer-facing information; ERP holds operational and financial
information. Connected, the salesperson never has to guess whether a customer
was delivered, invoiced, or paid.

\begin{erpfig}[The integrated customer view]{What a salesperson can see when CRM
and ERP share master data. The left column is relationship history; the right
column is operational truth that would otherwise require a phone call to another
department.}{fig:customer-360}
\begin{tikzpicture}[x=1cm,y=1cm]
\node[erpstore,minimum width=3.4cm,minimum height=1.3cm,draw=UTPNavy,fill=UTPNavy!8]
  (cust) at (7.3,0.9) {Seri Iskandar\\\scriptsize\bfseries Cycling Club};
\foreach \i/\t in {0/{Opportunities and pipeline stage},
                   1/{Campaigns and offers received},
                   2/{Service history and complaints}}{
  \pgfmathsetmacro{\y}{2.2-\i*1.3}
  \node[erpstep,minimum width=5.2cm,text width=4.8cm,minimum height=0.85cm] (l\i) at (2.7,\y) {\t};
  \draw[erpflow] (l\i) -- (cust);}
\foreach \i/\t in {0/{Open orders and delivery status},
                   1/{Invoices, payments, receivable balance},
                   2/{Purchase history, product mix, discounts}}{
  \pgfmathsetmacro{\y}{2.2-\i*1.3}
  \node[erpgood,minimum width=5.2cm,text width=4.8cm,minimum height=0.85cm] (r\i) at (11.9,\y) {\t};
  \draw[erpflow] (cust) -- (r\i);}
\node[erpcap,color=UTPBlue] at (2.7,3.25) {\bfseries from CRM};
\node[erpcap,color=UTPGreen] at (11.9,3.25) {\bfseries from ERP};
\node[erpnote,text width=5.0cm] at (7.3,-0.35)
  {one identity, so both halves describe the same customer};
\end{tikzpicture}
\end{erpfig}

This supports better decisions: which customers to prioritise, which products to
cross-sell, which accounts need service attention, and whether a discount is
commercially sensible.

\section{The Luna Electronics problem}

Tutorial 4 asks how CRM could help Luna Electronics grow its customer base,
retain existing customers, and improve service. A convincing answer connects
each business problem to a mechanism rather than asserting that ``CRM improves
customer relationships.''

\begin{workedexample}
Suppose Luna Electronics loses repeat customers because complaints are handled
through personal email inboxes. A CRM process logs the interaction against the
customer, assigns responsibility, stores the response, and makes the history
visible to the next employee. The benefit is not the database; it is faster,
more consistent service and less dependence on one person's memory.
\end{workedexample}

\begin{erptable}
\begin{tabularx}{\linewidth}{@{}>{\bfseries\leavevmode\color{ERPNavy}}p{0.23\linewidth}p{0.31\linewidth}X@{}}
\toprule
\rowcolor{ERPPaleBlue!92}
Business goal & CRM mechanism & Possible evidence of success \\
\midrule
Acquire customers & Lead capture, campaign tracking, opportunity pipeline, follow-up activities. & Lead-to-opportunity and opportunity-to-order conversion. \\
Retain customers & Interaction history, segmentation, targeted offers, proactive follow-up. & Repeat-purchase rate, churn rate, customer lifetime value trend. \\
Improve service & Shared case history, knowledge base, ownership and escalation. & Response time, resolution time, repeated-contact rate, satisfaction score. \\
Improve sales productivity & Pipeline visibility, next activities, opportunity value, automated document hand-off. & Sales cycle length, win rate, forecast accuracy. \\
\bottomrule
\end{tabularx}
\end{erptable}

\section{CRM data quality and governance}

CRM is only useful if staff record interactions consistently. The common failures
are duplicate customers, opportunities with no next activity, optimistic expected
revenue, inconsistent stages, and missing contact details.

The consequence is a pipeline that looks convincing and forecasts badly.

\begin{erpfig}[Discipline decides the forecast]{One pipeline, two readings. The
disciplined forecast weights each stage by a realistic win probability and lands
close to actual; reading the raw value as if every open deal will close
over-forecasts by more than double.}{fig:pipeline-quality}
\begin{tikzpicture}
\begin{axis}[
  width=13.4cm, height=5.6cm,
  ybar=2.5pt, bar width=9pt,
  ymin=0, ymax=430,
  ylabel={RM thousand},
  symbolic x coords={Qualified,Proposition,Negotiation,Won,Forecast,Actual},
  xtick=data, x tick label style={font=\scriptsize\color{UTPMuted}},
  enlarge x limits=0.11,
  legend style={at={(0.5,1.19)},anchor=north,legend columns=2},
  ymajorgrids, grid style={draw=UTPMuted!22},
  nodes near coords, every node near coord/.append style={font=\tiny\color{UTPMuted}},
]
\addplot[draw=UTPGreen,fill=UTPPaleTeal] coordinates
  {(Qualified,30) (Proposition,48) (Negotiation,60) (Won,45) (Forecast,183) (Actual,176)};
\addplot[draw=UTPRed,fill=UTPPaleRed] coordinates
  {(Qualified,100) (Proposition,120) (Negotiation,100) (Won,45) (Forecast,365) (Actual,176)};
\legend{weighted by stage probability (30\%{,} 40\%{,} 60\%{,} 100\%), raw pipeline value}
\end{axis}
\end{tikzpicture}
\end{erpfig}

A good CRM process therefore defines what counts as a lead and an opportunity,
when a stage should change, who owns the record, which fields are mandatory, how
duplicates are prevented, what ``won'' and ``lost'' mean, and how lost reasons
are recorded.

\begin{pitfall}
A colourful pipeline is not automatically a reliable sales forecast. If users
move opportunities between stages inconsistently or enter unrealistic expected
revenue, management gets an attractive but misleading report.
\end{pitfall}

\section{CRM in the West End Bicycles Odoo flow}

\begin{erpfig}[The Odoo pipeline]{The opportunity from the supplied guide moving
across the pipeline, and the point at which it becomes an operational document.
Expected revenue exists from the first stage; committed revenue does not exist
until the Sales Order is confirmed.}{fig:odoo-pipeline}
\begin{tikzpicture}[x=1cm,y=1cm]
\foreach \i/\st/\c in {0/New/UTPMuted, 1/Qualified/UTPBlue,
                       2/Proposition/UTPBlue, 3/Won/UTPGreen}{
  \pgfmathsetmacro{\x}{1.7+\i*3.6}
  \draw[draw=\c!55,fill=\c!8,rounded corners=2mm] (\x-1.55,0.0) rectangle (\x+1.55,2.55);
  \node[erpcap,color=\c] at (\x,2.28) {\bfseries \st};}
\node[erpdoc,minimum width=2.7cm,text width=2.4cm,minimum height=1.35cm] (o0) at (1.7,1.05)
  {Road Bike Order\\ --- Seri Iskandar};
\node[erpdoc,minimum width=2.7cm,text width=2.4cm,minimum height=1.35cm,draw=UTPBlue] (o1) at (5.3,1.05)
  {$+$ customer confirmed};
\node[erpdoc,minimum width=2.7cm,text width=2.4cm,minimum height=1.35cm,draw=UTPBlue] (o2) at (8.9,1.05)
  {$+$ quantity and price agreed};
\node[erpdoc,minimum width=2.7cm,text width=2.4cm,minimum height=1.35cm,draw=UTPGreen] (o3) at (12.5,1.05)
  {opportunity marked Won};
\foreach \i in {0,...,2}{\pgfmathtruncatemacro{\j}{\i+1}\draw[erpflow] (o\i)--(o\j);}
\node[erpchipwarn,anchor=west] at (0.15,-0.45) {expected revenue RM 5,000 --- a forecast, not income};
\node[erpstep,minimum width=3.4cm,minimum height=0.95cm,draw=UTPGreen,fill=UTPPaleTeal] (q) at (8.9,-1.85)
  {Quotation\\\tiny\mdseries 5 Road Bikes at RM 1,000};
\node[erpstep,minimum width=3.4cm,minimum height=0.95cm,draw=UTPGreen,fill=UTPPaleTeal] (so) at (13.2,-1.85)
  {Sales Order\\\tiny\mdseries RM 5,000 committed};
\draw[erpflow] (o2.south) -- node[erpnote,right,text width=2.4cm]{create quotation from the opportunity} (q.north);
\draw[erpflow] (q) -- node[erpnote,above]{confirm} (so);
\draw[erpflow] (so.north) -- (o3.south);
\end{tikzpicture}
\end{erpfig}

\begin{odoobox}
The supplied guide creates an opportunity called \textbf{Road Bike Order --- Seri
Iskandar Cycling Club} with expected revenue of RM 5,000, assigns the Sales
Manager, moves it through the pipeline, and then creates the quotation for 5 Road
Bikes. Three CRM principles are on display: the deal is owned by a named person,
expected revenue exists before the sale is confirmed, and the opportunity creates
the quotation so relationship history survives the hand-off.
\end{odoobox}

\section{How to evaluate a CRM proposal}

For the Tutorial 4 proposal, six questions make a strong evaluation framework:

\begin{enumerate}
  \item \textbf{Fit:} Does the system support the required sales, marketing, and service processes?
  \item \textbf{Integration:} Can customer, order, delivery, and financial data be connected?
  \item \textbf{Usability:} Will staff actually record and update interactions?
  \item \textbf{Reporting:} Can management see conversion, pipeline, campaigns, and service performance?
  \item \textbf{Governance and security:} Can access, ownership, and customer-data handling be controlled?
  \item \textbf{Cost and scalability:} Is the system affordable and able to grow with the business?
\end{enumerate}

\begin{tryit}
For Luna Electronics, choose one customer problem and write a chain in this
format: \emph{problem $\rightarrow$ CRM feature $\rightarrow$ changed employee
behaviour $\rightarrow$ measurable business outcome}. If one link is vague, the
proposal needs more detail.
\end{tryit}

\begin{quickcheck}
\begin{enumerate}
  \item Why is an opportunity different from a sales order?
  \item Give one example of cross-selling and one of up-selling for Luna Electronics.
  \item Why does CRM data quality affect sales forecasting?
\end{enumerate}
\end{quickcheck}

\begin{chapterrecap}
\begin{itemize}
  \item CRM answers what happens next and who owns it, which a contact list cannot (\figref{fig:contact-vs-crm}).
  \item The funnel narrows at every stage, so pipeline value is not revenue (\figref{fig:sales-funnel}).
  \item Integration with ERP adds delivery, invoice, and payment truth to relationship history (\figref{fig:customer-360}).
  \item Stage discipline, not software, decides whether a pipeline forecasts well (\figref{fig:pipeline-quality}).
  \item In Odoo the opportunity creates the quotation, so the history survives the hand-off to Sales (\figref{fig:odoo-pipeline}).
\end{itemize}
\end{chapterrecap}
